% Generated from TECHNICAL_REPORT.md. Do not edit this file directly.
% Run report/build_technical_report_latex.py after editing the Markdown source.

\FloatBarrier
\chapter{Executive summary}\label{chap:executive-summary}

This report describes a mobile-manipulation solution for the five JCIIOT 2026 RunningRobot factory-sorting levels. The robot must resolve a source and destination, navigate a large factory scene, grasp the specified container with two arms, transport it while avoiding stations and other materials, and physically release it at the target. L5 repeats the complete cycle for three neighboring totes while preserving already completed placements.

The solution is a geometry-grounded, fail-closed execution stack. A semantic plan selects \CodeInline{move \allowbreak{}-\allowbreak{}> \allowbreak{}pick\_\allowbreak{}up \allowbreak{}-\allowbreak{}> \allowbreak{}move \allowbreak{}-\allowbreak{}> \allowbreak{}place\_\allowbreak{}down}; an 8-connected A* planner produces collision-aware routes; a dual-arm operational-space controller performs approach, grasp, lift, lowering, and release; and a payload-aware departure planner prevents a carried tote from sweeping through adjacent source objects. The simulator state is recorded after every physical step. A separate realism auditor then checks frame-to-frame continuity, reconstructs MuJoCo contacts during every held frame, and rejects any recorded collision or unintended material contact.

\SystemArchitectureFigure

The evaluated L1-L5 runs score 10/10, 15/15, 20/20, 25/25, and 30/30, for 100/100 overall. They contain 18,332 recorded frames and seven successful grasps. Across all five runs, the maximum base translation is 0.035113 m per recorded frame, the maximum base rotation is 0.052434 rad per frame, recorded collision frames are zero, and held-material-to-unheld-material contact frames are zero. Aggregated by level, gripper contact is present for 99.5799\% to 100\% of held frames; the lowest per-object value in L5 is 99.3041\%. Each L1-L5 video contains exactly one rendered frame for every source trajectory frame and passes a complete decode test.

Navigation runs at 0.70 m/s and 20 Hz, limiting commanded translation to about 3.5 cm per frame. Efficiency comes from source-departure geometry, accurate staging, and replanning from open space - never from skipped physics frames, teleports, or direct object relocation.

\FloatBarrier
\chapter{Task, scope, and evidence model}\label{chap:task-scope-and-evidence-model}

\section{Task structure}\label{sec:task-structure}

L1-L4 each require one container transfer. L5 requires three transfers from the same crowded input station to the same auxiliary output table. The objective score is gated by successful grasp events, leaving the source region, arriving within the official target radius, and avoiding the judge collision penalty. The maximum scores are 10, 15, 20, 25, and 30 points.

Execution uses explicit official identifiers from \CodeInline{knowledge/\allowbreak{}task\_\allowbreak{}config.\allowbreak{}json}. This removes natural-language naming ambiguity while preserving the same task semantics. Each level is evaluated in an isolated process and environment. One fixed code, task, motion, scoring, and realism configuration is applied consistently across all five reported runs.

\section{Evidence hierarchy}\label{sec:evidence-hierarchy}

The deliverable uses four independent evidence layers:

\begin{enumerate}
  \item \CodeInline{trajectory.\allowbreak{}json} is the simulator-emitted state sequence, including robot joints, base pose, material poses, events, and collision flags.
  \item \CodeInline{score.\allowbreak{}json} applies scoring rule \CodeInline{grasp\_\allowbreak{}success\_\allowbreak{}gate\_\allowbreak{}l5\_\allowbreak{}multi\_\allowbreak{}v2} against official semantic-map centers at reference commit \CodeInline{129e94a9cff787031472045e19c24a4baeaefc48}.
  \item \CodeInline{realism\_\allowbreak{}audit.\allowbreak{}json} checks continuity and reconstructs contact pairs from each recorded state. It never edits a trajectory.
  \item \CodeInline{Lx\_\allowbreak{}dualview\_\allowbreak{}full.\allowbreak{}mp4} is rendered from the same trajectory. Its metadata binds the source trajectory and video with SHA-256 hashes; \CodeInline{video\_\allowbreak{}verification.\allowbreak{}json} records a complete decode of every frame.
\end{enumerate}

Score and realism are conjunctive release gates. A run is accepted only if it reaches the maximum objective score and passes every realism check. A nominal score alone is insufficient.

\section{Realism requirements}\label{sec:realism-requirements}

The implementation explicitly prohibits the following as runtime fallbacks:

\begin{itemize}
  \item grasp and release only through the physical skill path - no direct kinematic pick/place API, ungrasped placement, or one-step object relocation;
  \item no unrecorded base jump or carried-payload turn through neighboring materials;
  \item no ignored unheld-material contact, wrapped reset that differs from the last recorded navigation state, or fabricated missing-physics success.
\end{itemize}

\FloatBarrier
\chapter{Technical description}\label{chap:technical-description}

\section{Architecture}\label{sec:architecture}

The system separates task semantics, geometry, execution, and evidence:

\ArchitectureDetailFigure

Task-level decisions are deterministic. Geometric decisions use live simulator poses, not hard-coded destination teleports. Motion functions return failure if a required controller, attachment, path, clearance test, or contact condition is missing.

\section{Navigation and safe speed profile}\label{sec:navigation-and-safe-speed-profile}

\CodeInline{MoveSkill} resolves the named station to its official approach point and plans an 8-connected A* path on the generated occupancy grid. Obstacle inflation accounts for the mobile manipulator footprint. A clearance ladder allows the planner to relax margins only when a more conservative route is unavailable. Small start and goal neighborhoods are treated specially because valid station approach points lie close to station geometry.

The submitted base profile is:

\begin{table}[htbp]
  \centering
  \caption{Submitted navigation and safety parameters}
  \label{tab:navigation-profile}
  \begin{tblr}{
    width=\linewidth,
    colspec={Q[l,0.85] Q[c,0.58] Q[l,1.75]},
    row{1}={bg=ReportNavy,fg=white,font=\bfseries},
    row{even}={bg=ReportMist},
    cells={font=\footnotesize,valign=m},
    hlines={0.35pt,ReportRule},
    vlines={0.25pt,ReportRule},
    colsep=3.2pt,rowsep=3.2pt
  }
    Parameter & Configured value & Safety purpose \\
    Control frequency & 20 Hz & one bounded update per control frame \\
    Maximum linear speed & 0.70 m/s & approximately 0.035 m maximum commanded step \\
    Maximum angular speed & 1.20 rad/s & approximately 0.060 rad theoretical step \\
    Waypoint tolerance & 0.01 m & prevents coarse endpoint acceptance \\
    Endpoint settle & 5 steps & lets contacts and pose converge \\
    Transport retreat & 1.0 m & clears the source table before routing \\
    Transport lane & 2.2 m & clears neighboring source materials \\
  \end{tblr}
\end{table}

The backend advances the base in bounded increments and records every update. Each increment is followed by MuJoCo state propagation, attachment synchronization if an object is held, and collision inspection. The reported runs observe at most 0.035113 m translation and 0.052434 rad rotation per frame, both below the audit limits.

This base implementation is a kinematic simulator drive rather than wheel-torque control; that limitation is stated explicitly in Section 7. It is nevertheless temporally continuous at the submitted sampling rate and does not jump between distant poses.

\section{Geometry-derived dual-arm grasp}\label{sec:geometry-derived-dual-arm-grasp}

\CodeInline{PickUpSkill} first validates the requested material against the live scene and computes the approach from the object center and its model grasp sites. When a default site lies behind a scene proxy, a calibrated grasp frame is rotated to an accessible wall while preserving the actual object pose. The resulting base pose, yaw, and pinch points are passed to the scripted expert.

The scripted expert performs a physically stepped sequence:

\begin{enumerate}
  \item stow the arms before narrow navigation where required;
  \item move to a safe height with operational-space control;
  \item move horizontally above the live pinch sites;
  \item descend with closed-loop Cartesian feedback;
  \item close both grippers;
  \item validate finger-pad contact;
  \item lift the container and validate lifted height.
\end{enumerate}

The submitted grasp profile records every simulation step (\CodeInline{record\_\allowbreak{}frame\_\allowbreak{}interval \allowbreak{}= \allowbreak{}1}), allows up to 300 lift-control steps, uses a 0.02 m lift tolerance, and holds for 20 steps. The official collector's step function is wrapped so that intermediate controller steps are recorded, not merely phase endpoints.

If the physics grasp is unavailable or contact validation fails, the skill fails. There is no non-physics success fallback.

\section{Continuous handoff between navigation and grasp environments}\label{sec:continuous-handoff-between-navigation-and-grasp-environments}

Grasp execution uses a wrapped evaluation environment. A reset would normally restore all material objects to their spawn poses and could silently undo an already completed L5 placement. Before entering the wrapper, the solution snapshots the complete live material state and all non-base robot joints. The wrapper is initialized with that exact state, and its reset hook reapplies the snapshot before grasp motion begins.

After grasp, the target object's real grasp state is synchronized back to navigation while every non-target object retains its live pre-grasp state. This acts as a transactional state handoff:

\StateHandoffFigure

The procedure is especially important in L5, where the second and third wrapped grasps must preserve totes already delivered to the destination.

\section{Carried-object transport}\label{sec:carried-object-transport}

After dual-gripper contact and lift validation, the transport utility captures the object's transform relative to the robot. During base motion, that relative transform is updated at every recorded step. This is a continuous simulator attachment, not a source-to-target relocation. The realism audit independently confirms gripper-object contact for at least 99\% of held frames and checks that the held object does not contact any other material object.

Navigation collision logic ignores only the ground and the currently held object where necessary for the attachment. Fixtures, stations, and all unheld material objects remain collision-relevant. \CodeInline{follow\_\allowbreak{}path()} stops and returns failure when a forbidden contact is observed.

\section{Payload-aware source departure}\label{sec:payload-aware-source-departure}

The most demanding geometry occurs at the L5 input table: three target totes begin side by side, and rotating the front tote in place can sweep it through the center tote. The solution handles this geometry with a payload-aware departure plan:

\begin{enumerate}
  \item derive the outward direction from the held object to the base;
  \item select the tangential direction from the initial A* route, not simply the destination bearing;
  \item translate outward by 1.0 m;
  \item continue along a 2.2 m clearance lane;
  \item replan globally from the lane exit.
\end{enumerate}

The fine-grained passability check uses 0.05 m sampling. Only the first 0.10 m may use an endpoint escape mask so a conservative occupancy cell does not block a valid direct retreat; the remainder must be passable. If translation is impossible, an in-place payload turn is considered only after an arc preflight verifies at least 0.65 m clearance from every other material center. The check is fail-closed.

L5 processes the crowded source in \CodeInline{front \allowbreak{}-\allowbreak{}> \allowbreak{}center \allowbreak{}-\allowbreak{}> \allowbreak{}back} order. This exposes a safe outward corridor for each next object and avoids carrying one tote through another.

\section{Physical placement}\label{sec:physical-placement}

At the destination, the robot first retreats to an open staging point, turns continuously, and approaches the table in a straight line. Both arms and the attached object are lowered together using operational-space control. The grippers then open, the attachment is cleared, and MuJoCo is stepped for 60 release frames so the object can settle on the table.

The three L5 drops use lateral offsets of 0.00, +0.65, and -0.65 m relative to the target approach. Navigation and placement use the same offset, avoiding a sideways sweep at release. Terminal distances from the official target center are 0.065845 m for the front tote, 0.667618 m for the center tote, and 0.646361 m for the back tote, all inside the 0.80 m scoring radius.

\LFivePlacementFigure

\section{Per-step realism auditor}\label{sec:per-step-realism-auditor}

\CodeInline{audit\_\allowbreak{}trajectory\_\allowbreak{}realism.\allowbreak{}py} is separate from the executor and score calculator. It rejects a run if any threshold is exceeded:

\begin{table}[htbp]
  \centering
  \caption{Strict realism-audit acceptance limits}
  \label{tab:audit-limits}
  \begin{tblr}{
    width=\linewidth,
    colspec={Q[l,2.1] Q[c,0.75]},
    row{1}={bg=ReportNavy,fg=white,font=\bfseries},
    row{even}={bg=ReportMist},
    cells={font=\footnotesize,valign=m},
    hlines={0.35pt,ReportRule},
    vlines={0.25pt,ReportRule},
    colsep=3.2pt,rowsep=3.2pt
  }
    Quantity & Limit \\
    Base translation per frame & 0.06 m \\
    Base rotation per frame & 0.08 rad \\
    Robot joint change per frame & 0.35 \\
    Material translation per frame & 0.10 m \\
    Material rotation per frame & 0.25 rad \\
    Held-frame gripper contact fraction & at least 0.99 \\
    Recorded collision frames & 0 \\
    Held-material to other-material contact frames & 0 \\
  \end{tblr}
\end{table}

For each held frame, the auditor reconstructs the recorded state in the same MuJoCo scene, calls forward propagation, enumerates active contact pairs, and classifies contacts by robot, held material, unheld material, fixture, and ground. This catches a carried payload striking another tote even if a coarse base-only collision flag would miss it. Audit JSON records the worst transition, its frame number, missing-contact frames, unintended pairs, thresholds, and explicit failures.

\ContinuityMarginsFigure

\section{Full-frame video generation and verification}\label{sec:full-frame-video-generation-and-verification}

The demonstration renderer replays each submitted trajectory state by state. Each source frame produces exactly one H.264 frame at 20 FPS; \CodeInline{subsample\_\allowbreak{}step} is 1. The composite view is 768 x 288 pixels and pairs an environment overview with \CodeInline{robot0\_\allowbreak{}robotview}. L4 uses \CodeInline{agentview} for the overview because it gives a clearer scene composition; the other levels use \CodeInline{frontview}.

The video verifier opens each MP4 independently with OpenCV and decodes it to end-of-stream. It checks declared, decoded, metadata, and source frame counts; both view halves are checked for black frames; and SHA-256 hashes are recomputed. All five videos pass. Video is qualitative evidence and does not replace score or contact verification.

\VideoEvidenceFigure

\section{Implementation and third-party software}\label{sec:implementation-and-third-party-software}

The participant implementation is concentrated in the following files (paths are relative to \CodeInline{source\_\allowbreak{}code/\allowbreak{}}):

\begin{table}[htbp]
  \centering
  \caption{Participant implementation components}
  \label{tab:implementation}
  \begin{tblr}{
    width=\linewidth,
    colspec={Q[l,1.35] Q[l,2.35]},
    row{1}={bg=ReportNavy,fg=white,font=\bfseries},
    row{even}={bg=ReportMist},
    cells={font=\footnotesize,valign=m},
    hlines={0.35pt,ReportRule},
    vlines={0.25pt,ReportRule},
    colsep=3.2pt,rowsep=3.2pt
  }
    Path & Responsibility \\
    \CodeInline{robot\_\allowbreak{}agent/\allowbreak{}skills/\allowbreak{}move.\allowbreak{}py} & semantic target resolution, A*, L5 offsets, payload-aware departure \\
    \CodeInline{robot\_\allowbreak{}agent/\allowbreak{}skills/\allowbreak{}pick\_\allowbreak{}up.\allowbreak{}py} & live object validation, exact pose handoff, scene snapshot, physics-only grasp \\
    \CodeInline{robot\_\allowbreak{}agent/\allowbreak{}skills/\allowbreak{}scripted\_\allowbreak{}grasp.\allowbreak{}py} & per-step recording and closed-loop dual-arm expert \\
    \CodeInline{robot\_\allowbreak{}agent/\allowbreak{}skills/\allowbreak{}place\_\allowbreak{}down.\allowbreak{}py} & matched target offsets and physics-only placement \\
    \CodeInline{robot\_\allowbreak{}agent/\allowbreak{}environments/\allowbreak{}robosuite\_\allowbreak{}backend.\allowbreak{}py} & bounded base drive, collision checking, lowering, release, continuous turn \\
    \CodeInline{pipeline/\allowbreak{}audit\_\allowbreak{}trajectory\_\allowbreak{}realism.\allowbreak{}py} & offline continuity and reconstructed-contact audit \\
    \CodeInline{pipeline/\allowbreak{}render\_\allowbreak{}trajectory\_\allowbreak{}video.\allowbreak{}py} & full-frame dual-view video rendering \\
    \CodeInline{knowledge/\allowbreak{}robot\_\allowbreak{}params.\allowbreak{}json} & fixed motion and audit-relevant parameters \\
  \end{tblr}
\end{table}

Third-party and organizer software is acknowledged below. Versions are those used in the submitted environment.

\begin{table}[htbp]
  \centering
  \caption{Software dependencies, roles, and licenses}
  \label{tab:software}
  \begin{tblr}{
    width=\linewidth,
    colspec={Q[l,0.72] Q[c,0.52] Q[l,1.05] Q[l,1.65]},
    row{1}={bg=ReportNavy,fg=white,font=\bfseries},
    row{even}={bg=ReportMist},
    cells={font=\footnotesize,valign=m},
    hlines={0.35pt,ReportRule},
    vlines={0.25pt,ReportRule},
    colsep=3.2pt,rowsep=3.2pt
  }
    Component & Version & Use & License / source \\
    MuJoCo & 3.9.0 & dynamics, contacts, rendering & Apache-2.0; official computation documentation \cite{mujoco2026computation} \\
    robosuite fork & 1.5.2 & robot models, controllers, factory environments & MIT; robosuite paper \cite{zhu2020robosuite} \\
    robomimic & 0.5.0 & optional behavior-cloning infrastructure & MIT; robomimic paper \cite{mandlekar2021robomimic} \\
    NumPy & 1.26.4 & geometry and trajectory analysis & BSD-3-Clause \\
    SciPy & 1.15.3 & numerical utilities & BSD-3-Clause \\
    PyTorch & 2.7.0+cu126 & optional BC checkpoint runtime & BSD-style \\
    OpenCV Python & 4.8.1.78 & independent video decode verification & Apache-2.0 \\
    ReportLab & 5.0.0 & PDF report generation only & BSD-style \\
  \end{tblr}
\end{table}

The official scene assets, task definitions, and scoring reference are organizer-provided. The evaluated deterministic expert path does not require online LLM or network access.

\FloatBarrier
\chapter{Novelty statement and relation to prior work}\label{chap:novelty-statement-and-relation-to-prior-work}

\section{Claimed contribution}\label{sec:claimed-contribution}

The contribution is not a new universal robot-learning algorithm. It is a competition-specific reliability architecture for a known mobile-manipulation domain. Its distinctive elements are:

\begin{enumerate}
  \item \textbf{Score-plus-reality release gating.} Full objective score is necessary but not sufficient; continuity, contact reconstruction, collision state, and full-frame video integrity must also pass.
  \item \textbf{Payload-aware departure rather than base-only planning.} The first meters after grasp are planned using the carried object's geometry, route direction, neighboring material locations, and a swept-turn preflight.
  \item \textbf{Transactional cross-environment continuity.} A wrapped grasp rollout begins from the exact live scene and commits only the new target/robot state while preserving all non-target objects.
  \item \textbf{Physical fail-closed manipulation.} Missing grasp, lift, attachment, route, turn clearance, lowering, or release capability returns failure; direct pick/place and animation fallbacks are disabled.
  \item \textbf{Per-physics-step evidence.} The recorder instruments every physical controller step, enabling frame-level limits and reconstructed contact auditing rather than phase-level screenshots.
\end{enumerate}

Together these mechanisms address a common gap in fixed-scene competition systems: an objective evaluator may accept a final pose even when the intermediate motion is implausible. The proposed release gate makes the intermediate trajectory a first-class result.

\section{Relationship to established and current methods}\label{sec:relationship-to-established-and-current-methods}

A* is the classical foundation for minimum-cost heuristic graph search \cite{hart1968astar}. This solution retains A* and adds manipulator-footprint inflation, endpoint escape bounds, route-aligned source departure, and payload arc preflight. These additions are application-layer safety mechanisms, not a replacement for A*.

robosuite provides modular robot simulation and controller infrastructure \cite{zhu2020robosuite}, while robomimic studies strong offline imitation-learning baselines and the design factors that matter for robot manipulation \cite{mandlekar2021robomimic}. The submitted known-scene system uses a deterministic closed-loop expert because its decisions and intermediate states are directly auditable. DAgger formalizes how sequential imitation errors can compound when a learned policy encounters states outside its training distribution \cite{ross2011dagger}; our fixed-scene response is explicit geometry and fail-closed recovery rather than additional online data collection.

Recent manipulation systems pursue broader policy capability. Diffusion Policy models multimodal visuomotor actions with conditional diffusion \cite{chi2023diffusion}. Action Chunking with Transformers (ACT) targets fine-grained bimanual manipulation and predicts action sequences \cite{zhao2023act}. Mobile ALOHA combines whole-body teleoperation and imitation learning for bimanual mobile manipulation \cite{fu2025mobilealoha}. Those systems address data-driven generalization and complex behavior acquisition. This solution advances a different axis: deterministic execution, intermediate-state continuity, and auditability in five fixed official scenes.

\begin{table}[htbp]
  \centering
  \caption{Positioning relative to learned-policy methods}
  \label{tab:method-positioning}
  \begin{tblr}{
    width=\linewidth,
    colspec={Q[l,0.75] Q[l,1.55] Q[l,1.55]},
    row{1}={bg=ReportNavy,fg=white,font=\bfseries},
    row{even}={bg=ReportMist},
    cells={font=\footnotesize,valign=m},
    hlines={0.35pt,ReportRule},
    vlines={0.25pt,ReportRule},
    colsep=3.2pt,rowsep=3.2pt
  }
    Dimension & Diffusion / ACT / learned-policy direction & This solution \\
    Primary goal & learn expressive policies from demonstrations & reliable execution in known official geometry \\
    Generalization & central research objective & explicitly limited to L1-L5 scenes \\
    Data requirement & demonstrations and training & semantic maps, live geometry, optional checkpoint \\
    Intermediate safety evidence & policy and benchmark dependent & frame limits plus reconstructed contacts \\
    Failure behavior & policy dependent & fail closed at named skill boundaries \\
    Claimed result & broader task-policy capability & 100/100 plus five realism-audit passes \\
  \end{tblr}
\end{table}

\section{Novelty boundary}\label{sec:novelty-boundary}

We claim novelty for the integration and competition application of these mechanisms, not for A*, obstacle inflation, operational-space control, object attachments, or offset placement individually. The result does not demonstrate superiority over general-purpose learned policies, does not establish research-benchmark state of the art, and does not prove real-world safety. It demonstrates that all five fixed simulator levels can be solved at full objective score while retaining continuous, contact-audited intermediate motion.

\FloatBarrier
\chapter{Experimental protocol}\label{chap:experimental-protocol}

\section{Evaluation procedure}\label{sec:evaluation-procedure}

The same fixed code and \CodeInline{robot\_\allowbreak{}params.\allowbreak{}json} configuration is used for every reported run. Each level follows the same sequence:

\begin{enumerate}
  \item start a fresh headless MuJoCo process;
  \item execute the canonical official task;
  \item save the simulator-emitted \CodeInline{\_\allowbreak{}OK.\allowbreak{}json} trajectory without modification;
  \item compute the objective score against the pinned official reference;
  \item run the strict realism auditor;
  \item validate the run only if score and audit pass;
  \item render the validated trajectory at one video frame per trajectory frame;
  \item fully decode and hash-check the video.
\end{enumerate}

No trajectory JSON was edited between execution, scoring, auditing, rendering, and packaging. SHA-256 binds the video metadata to the trajectory used for rendering.

\section{Reproducibility controls}\label{sec:reproducibility-controls}

\begin{itemize}
  \item official reference commit and score-rule version are stored in each score and manifest;
  \item run stamps and environment identifiers are stored in \CodeInline{final\_\allowbreak{}run\_\allowbreak{}summary.\allowbreak{}json};
  \item every per-level evidence directory contains trajectory, score, realism audit, and result metadata;
  \item official submission ZIPs contain exactly \CodeInline{trajectory.\allowbreak{}json}, \CodeInline{score.\allowbreak{}json}, and \CodeInline{submission\_\allowbreak{}manifest.\allowbreak{}json};
  \item \CodeInline{SHA256SUMS.\allowbreak{}txt} covers the five submission ZIPs;
  \item the top-level bundle has its own checksum manifest;
  \item source overlays and the relevant robosuite factory-sorting overrides are included.
\end{itemize}

\FloatBarrier
\chapter{Results and analysis}\label{chap:results-and-analysis}

\section{Objective score and run size}\label{sec:objective-score-and-run-size}

\begin{table}[htbp]
  \centering
  \caption{Objective scores, trajectory sizes, and execution times}
  \label{tab:score-summary}
  \begin{tblr}{
    width=\linewidth,
    colspec={Q[c,0.42] Q[l,1.55] Q[c,0.50] Q[r,0.55] Q[r,0.72] Q[r,0.76]},
    row{1}={bg=ReportNavy,fg=white,font=\bfseries},
    row{even}={bg=ReportMist},
    cells={font=\footnotesize,valign=m},
    hlines={0.35pt,ReportRule},
    vlines={0.25pt,ReportRule},
    colsep=3.2pt,rowsep=3.2pt
  }
    Level & Environment & Score & Frames & Wall time (s) & Video duration (s) \\
    L1 & FactorySorting1\_3FO3ERFHISEM & 10/10 & 1,793 & 79.950 & 89.65 \\
    L2 & FactorySorting3\_3FO3ERRPH7X9 & 15/15 & 1,674 & 75.260 & 83.70 \\
    L3 & FactorySorting5\_3FO3ERTPXEUT & 20/20 & 2,402 & 101.243 & 120.10 \\
    L4 & FactorySorting7\_3FO3ERFKY9RN & 25/25 & 2,324 & 117.107 & 116.20 \\
    L5 & FactorySorting9\_3FO3ERT2C5FP & 30/30 & 10,139 & 548.334 & 506.95 \\
    \textbf{Total} & five environments & \textbf{100/100} & \textbf{18,332} & \textbf{921.894} & \textbf{916.60} \\
  \end{tblr}
\end{table}

Wall time is one measured end-to-end run per level on the same host, not an average. Video duration is exactly \CodeInline{frames \allowbreak{}/\allowbreak{} \allowbreak{}20 \allowbreak{}FPS}; it is not a claimed real-robot completion time. L5 dominates because it contains three complete grasp-transport-place cycles and records every physical control step.

\section{Continuity and contact results}\label{sec:continuity-and-contact-results}

\begin{table}[htbp]
  \centering
  \caption{Observed frame-to-frame continuity maxima}
  \label{tab:continuity}
  \begin{tblr}{
    width=\linewidth,
    colspec={Q[c,0.55] Q[r,0.82] Q[r,0.82] Q[r,0.78] Q[r,0.82]},
    row{1}={bg=ReportNavy,fg=white,font=\bfseries},
    row{even}={bg=ReportMist},
    cells={font=\footnotesize,valign=m},
    hlines={0.35pt,ReportRule},
    vlines={0.25pt,ReportRule},
    colsep=3.2pt,rowsep=3.2pt
  }
    Level & Base step (m) & Base turn (rad) & Joint step & Object step (m) \\
    L1 & 0.035100 & 0.026382 & 0.225176 & 0.048384 \\
    L2 & 0.035101 & 0.052271 & 0.225462 & 0.047969 \\
    L3 & 0.035113 & 0.026600 & 0.225186 & 0.035006 \\
    L4 & 0.035100 & 0.052434 & 0.299188 & 0.052564 \\
    L5 & 0.035113 & 0.027476 & 0.225283 & 0.086433 \\
    Audit limit & 0.060000 & 0.080000 & 0.350000 & 0.100000 \\
  \end{tblr}
\end{table}

\begin{table}[htbp]
  \centering
  \caption{Rotation, grasp-contact, and collision audit results}
  \label{tab:contact-audit}
  \begin{tblr}{
    width=\linewidth,
    colspec={Q[c,0.42] Q[r,0.72] Q[c,1.20] Q[c,0.68] Q[c,1.10] Q[c,0.48]},
    row{1}={bg=ReportNavy,fg=white,font=\bfseries},
    row{even}={bg=ReportMist},
    cells={font=\footnotesize,valign=m},
    hlines={0.35pt,ReportRule},
    vlines={0.25pt,ReportRule},
    colsep=3.2pt,rowsep=3.2pt
  }
    Level & Object turn (rad) & Held contact & Collision frames & Unintended material-contact frames & Audit \\
    L1 & 0.160527 & 954 / 957 = 99.6865\% & 0 & 0 & PASS \\
    L2 & 0.064961 & 1,050 / 1,050 = 100\% & 0 & 0 & PASS \\
    L3 & 0.041349 & 1,178 / 1,180 = 99.8305\% & 0 & 0 & PASS \\
    L4 & 0.110816 & 997 / 997 = 100\% & 0 & 0 & PASS \\
    L5 & 0.179854 & 4,267 / 4,285 = 99.5799\% & 0 & 0 & PASS \\
    Audit limit & 0.250000 & at least 99\% & 0 & 0 & — \\
  \end{tblr}
\end{table}

The largest L5 object translations occur during physical release and settling rather than source-to-target transport. They remain below 0.10 m per frame. Short missing gripper-contact intervals occur at attachment transitions or release boundaries; all levels remain above the 99\% threshold, and the object's motion remains continuous.

\section{L5 multi-object analysis}\label{sec:l5-multi-object-analysis}

\begin{table}[htbp]
  \centering
  \caption{L5 per-object contact and terminal-placement results}
  \label{tab:l5-objects}
  \begin{tblr}{
    width=\linewidth,
    colspec={Q[l,0.62] Q[r,0.72] Q[r,0.82] Q[r,0.88] Q[r,1.05]},
    row{1}={bg=ReportNavy,fg=white,font=\bfseries},
    row{even}={bg=ReportMist},
    cells={font=\footnotesize,valign=m},
    hlines={0.35pt,ReportRule},
    vlines={0.25pt,ReportRule},
    colsep=3.2pt,rowsep=3.2pt
  }
    Object & Held frames & Contact frames & Contact fraction & Terminal target distance \\
    front & 1,411 & 1,403 & 99.4330\% & 0.065845 m \\
    center & 1,437 & 1,427 & 99.3041\% & 0.667618 m \\
    back & 1,437 & 1,437 & 100\% & 0.646361 m \\
  \end{tblr}
\end{table}

All three totes are grasped, moved out of the source region, and placed inside the target radius. Completed placements remain at the destination while subsequent wrapped grasps execute. The zero unintended-contact count confirms that none of the three carried totes is transported through another material object in the evaluated run.

\section{Video results}\label{sec:video-results}

\begin{table}[htbp]
  \centering
  \caption{Full-frame video verification results}
  \label{tab:video-results}
  \begin{tblr}{
    width=\linewidth,
    colspec={Q[c,0.38] Q[l,1.42] Q[r,0.72] Q[c,0.38] Q[c,0.72] Q[c,0.58] Q[c,0.48]},
    row{1}={bg=ReportNavy,fg=white,font=\bfseries},
    row{even}={bg=ReportMist},
    cells={font=\footnotesize,valign=m},
    hlines={0.35pt,ReportRule},
    vlines={0.25pt,ReportRule},
    colsep=3.2pt,rowsep=3.2pt
  }
    Level & File & Frames decoded & FPS & Resolution & Black frames & Result \\
    L1 & \CodeInline{L1\_\allowbreak{}dualview\_\allowbreak{}full.\allowbreak{}mp4} & 1,793 & 20 & 768 x 288 & 0 & PASS \\
    L2 & \CodeInline{L2\_\allowbreak{}dualview\_\allowbreak{}full.\allowbreak{}mp4} & 1,674 & 20 & 768 x 288 & 0 & PASS \\
    L3 & \CodeInline{L3\_\allowbreak{}dualview\_\allowbreak{}full.\allowbreak{}mp4} & 2,402 & 20 & 768 x 288 & 0 & PASS \\
    L4 & \CodeInline{L4\_\allowbreak{}dualview\_\allowbreak{}full.\allowbreak{}mp4} & 2,324 & 20 & 768 x 288 & 0 & PASS \\
    L5 & \CodeInline{L5\_\allowbreak{}dualview\_\allowbreak{}full.\allowbreak{}mp4} & 10,139 & 20 & 768 x 288 & 0 & PASS \\
  \end{tblr}
\end{table}

The decoded counts exactly equal the trajectory and metadata counts. Neither half of any composite video contains a black frame. The videos show each complete evaluated trajectory with one-to-one frame coverage rather than highlight-only excerpts.

\section{Strengths}\label{sec:strengths}

\begin{itemize}
  \item perfect objective score on all five official scenes;
  \item explicit prevention of direct pick/place and pose-jump fallbacks;
  \item per-step recording across navigation and controller internals;
  \item payload-aware departure and collision checks that include other materials;
  \item transparent, machine-readable realism limits and worst-case values;
  \item exact L1-L5 video/trajectory frame correspondence;
  \item deterministic execution without network services.
\end{itemize}

\section{Speed analysis and safe opportunities}\label{sec:speed-analysis-and-safe-opportunities}

The submitted profile prioritizes credible motion over minimum wall time. Runtime efficiency comes from clearing the source table through a route-derived departure lane and replanning from open space. Exact grasp-state handoff preserves completed L5 placements and prevents redundant manipulation cycles.

Additional safe optimization should target computation rather than larger physical jumps:

\begin{itemize}
  \item cache static collision geometry used by reconstructed-contact audits;
  \item simplify A* paths while retaining the 0.06 m frame limit and continuous collision tests;
  \item render videos in parallel after trajectories are recorded and validated;
  \item profile controller settling and reduce only frames with established convergence;
  \item encode videos asynchronously, since encoding does not affect robot physics.
\end{itemize}

Raising base steps toward the audit limit or reducing contact settling could shorten runs, but it would reduce safety margin and would require a complete L1-L5 evaluation and audit before adoption. The reported configuration retains the stated margins.

\FloatBarrier
\chapter{Limitations}\label{chap:limitations}

\begin{enumerate}
  \item \textbf{Fixed-scene evaluation.} Results cover five known official simulator scenes. The method is not evaluated on unseen layouts, object shapes, lighting, or randomized dynamics.
  \item \textbf{No real-robot deployment.} MuJoCo contact realism does not prove hardware safety. Sensor noise, actuator delay, compliance, wheel slip, calibration error, and emergency-stop design are outside this evaluation.
  \item \textbf{Kinematic mobile base.} The base is advanced through bounded simulator joint increments rather than wheel-torque dynamics. The trajectory is temporally continuous, but it is not a wheel-controller validation.
  \item \textbf{Simulator attachment while carrying.} A gripper-relative transport attachment maintains the grasp during mobile motion. Its per-frame motion and contact are audited, but it is not equivalent to modeling all real grasp forces.
  \item \textbf{Single evaluated run per level.} The report presents one deterministic run per level, not repeated-trial success rates or confidence intervals.
  \item \textbf{Engineering thresholds.} Continuity and contact limits are conservative competition guardrails, not formal safety certificates.
  \item \textbf{Video presentation.} The overview camera is intentionally wide to show the factory; small objects can be difficult to see there, so the paired robot camera is needed for manipulation detail.
\end{enumerate}

These limitations bound the novelty and performance claims. The evidence supports full-score, continuity-audited execution in the provided simulation, not universal physical realism.

\FloatBarrier
\chapter{Reproduction and package guide}\label{chap:reproduction-and-package-guide}

The submission bundle layout is:

\begin{ReportCode}
README.md
TECHNICAL_REPORT.md
output/pdf/TECHNICAL_REPORT.pdf
submission/L1_*.zip ... L5_*.zip
submission/SHA256SUMS.txt
evidence/L1 ... L5/
  trajectory.json
  score.json
  realism_audit.json
  run_result.json
videos/L1_dualview_full.mp4 ... L5_dualview_full.mp4
videos/*.metadata.json
results/final_run_summary.json
results/video_verification.json
results/package_verification.json
source_code/
verify_submission.py
FINAL_CHECKSUMS.sha256
\end{ReportCode}

Run \CodeInline{python \allowbreak{}verify\_\allowbreak{}submission.\allowbreak{}py} from the extracted bundle to recompute official-package integrity and scores. The packaged verification report records the expected PASS result. \CodeInline{results/\allowbreak{}final\_\allowbreak{}run\_\allowbreak{}summary.\allowbreak{}json} is the compact source for every quantitative table in this report.

\FloatBarrier
\chapter{Conclusion}\label{chap:conclusion}

The JCIIOT RunningRobot solution achieves 100/100 across L1-L5 while preserving continuous intermediate motion and explicit contact evidence. Its key contribution is a strict evidence gate: physical grasp and release, bounded base motion, payload-aware departure, transactional state handoff, reconstructed-contact auditing, and full-frame video verification all have to agree before an artifact is accepted.

This approach trades some execution time for stronger physical credibility and auditability. The submitted evidence comprises five score-maximizing packages, five PASS realism audits, and five source-complete demonstrations.

\FloatBarrier
\chapter{References}\label{chap:references}
\printbibliography[heading=none]
